\documentclass[10pt,a4paper]{amsart}
\usepackage{amsmath,amssymb,amsthm,graphicx}
\usepackage{verbatim}
\newenvironment{reference}{\comment}{\endcomment}
%\newenvironment{reference}{}{}
\newenvironment{slogan}{\comment}{\endcomment}
\newenvironment{history}{\comment}{\endcomment}

% For commutative diagrams we use Xy-pic
\usepackage[all]{xy}

% We use 2cell for 2-commutative diagrams.
\xyoption{2cell}
% Theorem environments.
%
\theoremstyle{plain}
\newtheorem{theorem}[subsection]{Theorem}
\newtheorem{proposition}[subsection]{Proposition}
\newtheorem{lemma}[subsection]{Lemma}

\theoremstyle{definition}
\newtheorem{definition}[subsection]{Definition}
\newtheorem{example}[subsection]{Example}
\newtheorem{exercise}[subsection]{Exercise}
\newtheorem{situation}[subsection]{Situation}

\theoremstyle{remark}
\newtheorem{remark}[subsection]{Remark}
\newtheorem{remarks}[subsection]{Remarks}

\numberwithin{equation}{subsection}

% Macros
%
\def\lim{\mathop{\mathrm{lim}}\nolimits}
\def\colim{\mathop{\mathrm{colim}}\nolimits}
\def\Spec{\mathop{\mathrm{Spec}}}
\def\Hom{\mathop{\mathrm{Hom}}\nolimits}
\def\Ext{\mathop{\mathrm{Ext}}\nolimits}
\def\SheafHom{\mathop{\mathcal{H}\!\mathit{om}}\nolimits}
\def\SheafExt{\mathop{\mathcal{E}\!\mathit{xt}}\nolimits}
\def\Sch{\mathit{Sch}}
\def\Mor{\mathop{\mathrm{Mor}}\nolimits}
\def\Ob{\mathop{\mathrm{Ob}}\nolimits}
\def\Sh{\mathop{\mathit{Sh}}\nolimits}
\def\NL{\mathop{N\!L}\nolimits}
\def\CH{\mathop{\mathrm{CH}}\nolimits}
\def\proetale{{pro\text{-}\acute{e}tale}}
\def\etale{{\acute{e}tale}}
\def\QCoh{\mathit{QCoh}}
\def\Ker{\mathop{\mathrm{Ker}}}
\def\Im{\mathop{\mathrm{Im}}}
\def\Coker{\mathop{\mathrm{Coker}}}
\def\Coim{\mathop{\mathrm{Coim}}}

% Boxtimes
%
\DeclareMathSymbol{\boxtimes}{\mathbin}{AMSa}{"02}

%
% Macros for moduli stacks/spaces
%
\def\QCohstack{\mathcal{QC}\!\mathit{oh}}
\def\Cohstack{\mathcal{C}\!\mathit{oh}}
\def\Spacesstack{\mathcal{S}\!\mathit{paces}}
\def\Quotfunctor{\mathrm{Quot}}
\def\Hilbfunctor{\mathrm{Hilb}}
\def\Curvesstack{\mathcal{C}\!\mathit{urves}}
\def\Polarizedstack{\mathcal{P}\!\mathit{olarized}}
\def\Complexesstack{\mathcal{C}\!\mathit{omplexes}}
% \Pic is the operator that assigns to X its picard group, usage \Pic(X)
% \Picardstack_{X/B} denotes the Picard stack of X over B
% \Picardfunctor_{X/B} denotes the Picard functor of X over B
\def\Pic{\mathop{\mathrm{Pic}}\nolimits}
\def\Picardstack{\mathcal{P}\!\mathit{ic}}
\def\Picardfunctor{\mathrm{Pic}}
\def\Deformationcategory{\mathcal{D}\!\mathit{ef}}
\usepackage{cleveref}
\crefname{originallemma}{original lemma}{original lemmas}
\Crefname{originallemma}{Original lemma}{Original lemmas}
\makeatletter
\newlabel{decent-spaces-lemma-decent-specialization}{{03IL}{source}{}{sourcepointer-1}{}}
\newlabel{decent-spaces-lemma-decent-specialization@cref}{{[originallemma][0][]03IL}{[1][0][]source}{}{}{}}
\makeatother
\crefname{originallemma}{original lemma}{original lemmas}
\Crefname{originallemma}{Original lemma}{Original lemmas}
\makeatletter
\newlabel{decent-spaces-lemma-topology-fibre}{{0AC8}{source}{}{sourcepointer-2}{}}
\newlabel{decent-spaces-lemma-topology-fibre@cref}{{[originallemma][0][]0AC8}{[1][0][]source}{}{}{}}
\makeatother
\crefname{originallemma}{original lemma}{original lemmas}
\Crefname{originallemma}{Original lemma}{Original lemmas}
\makeatletter
\newlabel{flat-lemma-associated-point-specializes}{{05L3}{source}{}{sourcepointer-3}{}}
\newlabel{flat-lemma-associated-point-specializes@cref}{{[originallemma][0][]05L3}{[1][0][]source}{}{}{}}
\makeatother
\title{The closure of an associated point that meets the closed fibre meets an associated point of the fibre}
\author{The Stacks Project Authors}
\date{stacks-project:spaces-flat.tex:local-snapshot-20261012; GNU FDL 1.2 or later}
\begin{document}
\maketitle
\noindent\textit{Editorial note.} All delivered mathematical text is the author's own, copied from the pinned Stacks Project chapter file spaces-flat.tex; no statement or proof step was rewritten, reordered or removed.\par\medskip
\noindent\textit{Editorial note.} The counted proof is the author's complete argument: a lift of the specialisation along an etale chart of a decent space, an identification of the relative assassin through fibre associated points, and an application of the schemes case\par\medskip
\noindent\textit{Editorial note.} The label lemma-associated-point-specializes names the counted statement, and the proof environment immediately following it in the source is delivered verbatim as its proof.\par\medskip
\noindent\textit{Editorial note.} The counted statement is carried by the author's own lemma environment in the chapter source.\par\medskip
\noindent\textit{Allowed standard prerequisites (editorial).}\begin{itemize}
\item The relative assassin of a finite type module over the base, and its description as the set of points whose fibre point is associated to the fibre of the module
\item Etale charts of a decent space and the lifting of specialisations from the space to the chart
\item The schemes version of the statement, identifying the associated points of the special fibre of the pullback with the fibre of the relative assassin
\end{itemize}\medskip
\ifdefined\hypertarget\hypertarget{sourcepointer-1}{}\fi
\noindent\textit{Source reference.} Original lemma 03IL: sources/stacks-chow/tags.txt, source lines 4607--4607. Prerequisite pointer invoked at proof line(s) 1015; it names a standard fact of a sibling chapter that the counted argument uses, and every step of the counted proof built on it is retained verbatim in the delivered body.\par\medskip
\ifdefined\hypertarget\hypertarget{sourcepointer-2}{}\fi
\noindent\textit{Source reference.} Original lemma 0AC8: sources/stacks-chow/tags.txt, source lines 13046--13046. Prerequisite pointer invoked at proof line(s) 1011; it names a standard fact of a sibling chapter that the counted argument uses, and every step of the counted proof built on it is retained verbatim in the delivered body.\par\medskip
\ifdefined\hypertarget\hypertarget{sourcepointer-3}{}\fi
\noindent\textit{Source reference.} Original lemma 05L3: sources/stacks-chow/tags.txt, source lines 7222--7222. Prerequisite pointer invoked at proof line(s) 1022; it names a standard fact of a sibling chapter that the counted argument uses, and every step of the counted proof built on it is retained verbatim in the delivered body.\par\medskip
\begin{lemma}
\label{lemma-associated-point-specializes}
Let $S$ be a scheme. Let $X$ be a decent algebraic space
locally of finite type over $S$.
Let $\mathcal{F}$ be a finite type, quasi-coherent $\mathcal{O}_X$-module.
Let $s \in S$ such that $\mathcal{F}$ is flat over $S$ at all points of $X_s$.
Let $x' \in \text{Ass}_{X/S}(\mathcal{F})$. If
the closure of $\{x'\}$ in $|X|$ meets $|X_s|$, then the closure
meets $\text{Ass}_{X/S}(\mathcal{F}) \cap |X_s|$.
\end{lemma}

\begin{proof}
Observe that $|X_s| \subset |X|$ is the set of points of $|X|$
lying over $s \in S$, see
Decent Spaces, Lemma \ref{decent-spaces-lemma-topology-fibre}.
Let $t \in |X_s|$ be a specialization of $x'$ in $|X|$.
Choose an affine scheme $U$ and a point $u \in U$ and
an \'etale morphism $\varphi : U \to X$ mapping $u$ to $t$.
By Decent Spaces, Lemma \ref{decent-spaces-lemma-decent-specialization}
we can choose a specialization $u' \leadsto u$
with $u'$ mapping to $x'$. Set $g = f \circ \varphi$.
Observe that $s' = g(u') = f(x')$ specializes to $s$.
By our definition of $\text{Ass}_{X/S}(\mathcal{F})$
we have $u' \in \text{Ass}_{U/S}(\varphi^*\mathcal{F})$.
By the schemes version of this lemma
(More on Flatness, Lemma \ref{flat-lemma-associated-point-specializes})
we see that there is a specialization $u' \leadsto u$ with
$u \in \text{Ass}_{U_s}(\varphi^*\mathcal{F}_s) =
\text{Ass}_{U/S}(\varphi^*\mathcal{F}) \cap U_s$.
Hence $x = \varphi(u) \in \text{Ass}_{X/S}(\mathcal{F})$
lies over $s$ and the lemma is proved.
\end{proof}

\end{document}
